% Part: model-theory
% Chapter: models-of-arithmetic
% Section: standard-models

\documentclass[../../../include/open-logic-section]{subfiles}

\begin{document}

\olfileid{mod}{mar}{stm}
\olsection{అంకగణితపు ప్రామాణిక నమూనాలు}

అంకగణిత భాష $\Lang{L_A}$ సంఖ్యలను, మరింత నిర్దిష్టంగా సహజ
సంఖ్యలను ఉద్దేశించినదని స్పష్టమే. అందువల్ల “ఆ” ప్రామాణిక
నమూనా~$\Struct{N}$ ప్రత్యేకమైనది: మనం చర్చించదలచిన నమూనా అదే.
అయితే తర్కశాస్త్రంలో మనకు తరచుగా నిర్మాణాత్మక ధర్మాలపైనే ఆసక్తి;
సమరూపమైన ఏ రెండు \tetoken{నిర్మాణాలకైనా}{!!{structure}s} అవి ఒకటే.
కాబట్టి కొంచెం విశాలంగా, $\Struct{N}$తో సమరూపమైన ప్రతి
\tetoken{నిర్మాణాన్నీ}{!!{structure}} “ప్రామాణికం”గా పరిగణించవచ్చు.

\begin{defn}
$\Lang{L_A}$కు గల ఒక \tetoken{నిర్మాణం}{!!{structure}},
$\Struct{N}$తో సమరూపమైతే \emph{ప్రామాణికం}.
\end{defn}

\begin{prop}
\ollabel{prop:standard-domain} ఒక
\tetoken{నిర్మాణం}{!!a{structure}}~$\Struct{M}$ ప్రామాణికమైతే,
దాని వ్యక్తి క్షేత్రం ప్రామాణిక సంఖ్యాపదాల విలువల సమితి, అంటే
\[
\Domain{M} = \Setabs{\Value{\num{n}}{M}}{n \in \Nat}
\]
\end{prop}

\begin{proof}
ప్రతి $\Value{\num{n}}{M}\in\Domain{M}$ అని స్పష్టమే. ప్రతి
$x\in\Domain{M}$ ఏదో $n$కు $\Value{\num{n}}{M}$తో సమానమని మాత్రమే
చూపాలి. $\Struct{M}$ ప్రామాణికం కాబట్టి అది $\Struct{N}$తో
సమరూపం. $g\colon\Nat\to\Domain{M}$ ఒక సమరూపత అనుకుందాం. అప్పుడు
$g(n)=g(\Value{\num{n}}{N})=\Value{\num{n}}{M}$. అంతేకాక, $g$
\tetoken{సంగ్రస్త}{!!{surjective}} కాబట్టి ప్రతి
$x\in\Domain{M}$కూ $g(n)=x$ అయ్యే $n\in\Nat$ ఉంటుంది.
\end{proof}

\begin{explain}
$\Lang{L_A}$కు గల నిర్మాణం~$\Struct{M}$ ప్రామాణికమైతే, దాని
\tetoken{వ్యక్తి క్షేత్రంలోని}{!!{domain}} మూలకాలన్నింటినీ
ప్రామాణిక సంఖ్యాపదాలు $\num{0}$, $\num{1}$, $\num{2}$, \dots,
అంటే $\Obj{0}$, $\Obj{0}'$, $\Obj{0}''$ మొదలైన పదాలు
పేరు పెట్టగలవు. దీని అర్థం $\Domain{M}$లోని
\tetoken{మూలకాలు}{!!{element}s} స్వయంగా సంఖ్యలే అని కాదు;
$\Domain{N}$లోని సంఖ్యలను ఎంచుకున్నట్లే వీటినీ ఎంచగలమని మాత్రమే.
\end{explain}

\begin{prob}
\olref[mod][mar][stm]{prop:standard-domain} యొక్క విపర్యయం అసత్యమని
చూపండి; అంటే
$\Domain{M}=\Setabs{\Value{\num{n}}{M}}{n\in\Nat}$ అయినా
$\Struct{N}$తో సమరూపం కాని ఒక
\tetoken{నిర్మాణం}{!!a{structure}}~$\Struct{M}$కు ఉదాహరణ ఇవ్వండి.
\end{prob}

\begin{prop}
\ollabel{prop:thq-standard}
$\Sat{M}{\Th{Q}}$ మరియు
$\Domain{M}=\Setabs{\Value{\num{n}}{M}}{n\in\Nat}$ అయితే,
$\Struct{M}$ ప్రామాణికం.
\end{prop}

\begin{proof}
$\Struct{M}$, $\Struct{N}$తో సమరూపమని చూపాలి.
$g(n)=\Value{\num{n}}{M}$గా నిర్వచించిన
$g\colon\Nat\to\Domain{M}$ను పరిశీలించండి. పరికల్పన వల్ల $g$
\tetoken{సంగ్రస్త}{!!{surjective}}. అది
\tetoken{ఒకటి-ఒకటి}{!!{injective}} కూడా: $n\neq m$ అయినప్పుడల్లా
$\Th{Q}\Proves\eq/[\num{n}][\num{m}]$. అందువల్ల
$\Sat{M}{\Th{Q}}$ కనుక $n\neq m$ అయినప్పుడల్లా
$\Sat{M}{\eq/[\num{n}][\num{m}]}$. కనుక $n\neq m$ అయితే
$\Value{\num{n}}{M}\neq\Value{\num{m}}{M}$, అంటే $g(n)\neq g(m)$.

$g$ ఒక సమరూపత అని కూడా నిర్ధారించాలి.
\begin{enumerate}
\item $\Assign{\Obj{0}}{N}=0$ కాబట్టి
  $g(\Assign{\Obj{0}}{N})=g(0)$. $g$ నిర్వచనం వల్ల
  $g(0)=\Value{\num{0}}{M}$. కానీ $\num{0}$ అనేది $\Obj{0}$;
  యాదృచ్ఛికంగా ఒక \tetoken{స్థిరసంకేతమైన}{!!a{constant}} పదం
  విలువను, ఆ \tetoken{నిర్మాణం}{!!{structure}} ఆ
  \tetoken{స్థిరసంకేతానికి}{!!{constant}} కేటాయించినదే ఇస్తుంది:
  $\Value{\Obj{0}}{M}=\Assign{\Obj{0}}{M}$. కాబట్టి అవసరమైనట్లుగా
  $g(\Assign{\Obj{0}}{N})=\Assign{\Obj{0}}{M}$.
\item $\Struct{N}$లో $\prime$ అనేది $\Nat$పై ఉత్తరాధికారి ప్రమేయం
  కాబట్టి $g(\Assign{\prime}{N}(n))=g(n+1)$. $g$ నిర్వచనం వల్ల
  $g(n+1)=\Value{\num{n+1}}{M}$. కానీ $\num{n+1}$, $\num{n}'$ ఒకే
  పదం; కనుక $\Value{\num{n+1}}{M}=\Value{\num{n}'}{M}$. విలువ
  ప్రమేయ నిర్వచనం వల్ల ఇది
  $=\Assign{\prime}{M}(\Value{\num{n}}{M})$. మరి
  $\Value{\num{n}}{M}=g(n)$ కాబట్టి
  $g(\Assign{\prime}{N}(n))=\Assign{\prime}{M}(g(n))$.
\item $\Struct{N}$లో $+$ అనేది $\Nat$పై సంకలన ప్రమేయం కాబట్టి
  $g(\Assign{+}{N}(n,m))=g(n+m)$. $g$ నిర్వచనం వల్ల
  $g(n+m)=\Value{\num{n+m}}{M}$. అయితే
  $\Th{Q}\Proves\num{n+m}=(\num{n}+\num{m})$; కనుక
  $\Value{\num{n+m}}{M}=\Value{\num{n}+\num{m}}{M}$.
  విలువ ప్రమేయ నిర్వచనం వల్ల ఇది
  $=\Assign{+}{M}(\Value{\num{n}}{M},\Value{\num{m}}{M})$.
  $\Value{\num{n}}{M}=g(n)$, $\Value{\num{m}}{M}=g(m)$ కాబట్టి
  $g(\Assign{+}{N}(n,m))=\Assign{+}{M}(g(n),g(m))$.
\item $g(\Assign{\times}{N}(n,m))
  =\Assign{\times}{M}(g(n),g(m))$: అభ్యాసం.
\item $\tuple{n,m}\in\Assign{<}{N}$ అవడం సరిగ్గా $n<m$ అయినప్పుడే.
  $n<m$ అయితే $\Th{Q}\Proves\num{n}<\num{m}$, అలాగే
  $\Sat{M}{\num{n}<\num{m}}$. కనుక
  $\tuple{\Value{\num{n}}{M},\Value{\num{m}}{M}}\in\Assign{<}{M}$,
  అంటే $\tuple{g(n),g(m)}\in\Assign{<}{M}$. $n\not<m$ అయితే
  $\Th{Q}\Proves\lnot\num{n}<\num{m}$, అందువల్ల
  $\Sat/{M}{\num{n}<\num{m}}$. కనుక మునుపటిలాగే
  $\tuple{g(n),g(m)}\notin\Assign{<}{M}$. రెండింటినీ కలిపితే,
  $\tuple{n,m}\in\Assign{<}{N}$ అవడం సరిగ్గా
  $\tuple{g(n),g(m)}\in\Assign{<}{M}$ అయినప్పుడే.
\end{enumerate}
\end{proof}

\begin{explain}
$\Lang{L_A}$కు గల మరే
\tetoken{నిర్మాణం}{!!{structure}}~$\Struct{M}$ అయినా
$\num{0}$, $\num{1}$, $\num{2}$ మొదలైనవి పేరు పెట్టే
\tetoken{మూలకాలను}{!!{element}s} కలిగి ఉంటుంది. అందువల్ల
$\Nat$ నుంచి దాని వ్యక్తి క్షేత్రానికి చిత్రణను నిర్వచించడానికి
$g$ అత్యంత సహజమైన మార్గం. కాబట్టి $\Struct{M}$ కనీసం
$\num{n}$ల గురించిన కొన్ని మౌలిక వాక్యాలను $\Struct{N}$లాగే సత్యం
చేస్తూ, $g$ ద్విగుణమైతే, $g$ సమరూపత అవడంలో ఆశ్చర్యం లేదు. నిజానికి,
$\Domain{M}$లో $\num{n}$లు పేరు పెట్టేవి తప్ప మరే
\tetoken{మూలకాలు}{!!{element}s} లేకపోతే, అది ఒక్కటే సమరూపత.
\end{explain}

\begin{prop}
\ollabel{prop:thq-unique-iso} $\Struct{M}$ ప్రామాణికమైతే,
\olref{prop:thq-standard} నిరూపణలోని $g$ మాత్రమే
$\Struct{N}$ నుంచి $\Struct{M}$కు గల సమరూపత.
\end{prop}

\begin{proof}
$h\colon\Nat\to\Domain{M}$ అనేది $\Struct{N}$, $\Struct{M}$ల మధ్య
ఒక సమరూపత అనుకుందాం. $n$పై ఆగమనం ద్వారా $g=h$ అని చూపుతాం.
$n=0$ అయితే, $g$ నిర్వచనం వల్ల
$g(0)=\Assign{\Obj{0}}{M}$. కానీ $h$ సమరూపత కాబట్టి
$h(0)=h(\Assign{\Obj{0}}{N})=\Assign{\Obj{0}}{M}$; కనుక
$g(0)=h(0)$.

ఇప్పుడు $n+1$ సందర్భాన్ని పరిశీలించండి. మనకు
\begin{align*}
  g(n+1) & = \Value{\num{n+1}}{M} \text{ $g$ నిర్వచనం వల్ల}\\
  & = \Value{\num{n}'}{M} \text{ ఎందుకంటే $\num{n+1}\ident\num{n}'$}\\
  & = \Assign{\prime}{M}(\Value{\num{n}}{M})
    \text{ $\Value{t'}{M}$ నిర్వచనం వల్ల}\\
  & = \Assign{\prime}{M}(g(n)) \text{ $g$ నిర్వచనం వల్ల}\\
  & = \Assign{\prime}{M}(h(n)) \text{ ఆగమన పరికల్పన వల్ల}\\
  & = h(\Assign{\prime}{N}(n)) \text{ $h$ సమరూపత కాబట్టి}\\
  & = h(n+1)
\end{align*}
\end{proof}

\begin{explain}
ప్రతి \tetoken{అనంతంగా లెక్కించదగిన}{!!{denumerable}} సమితి~$M$ను
తీసుకుంటే, $\Nat$కూ $M$కూ మధ్య ఒక \tetoken{ద్విగుణ ప్రమేయం}{!!a{bijection}}
ఉంటుంది. కాబట్టి అటువంటి ప్రతి $M$, ఒక ప్రామాణిక నమూనా
$\Struct{M}$కు సంభావ్య \tetoken{వ్యక్తి క్షేత్రం}{!!{domain}}.
నిజానికి, $z\in M$ను, తగిన $s$ ప్రమేయాన్ని వరుసగా
$\Assign{\Obj{0}}{M}$, $\Assign{\prime}{M}$గా ఎంచుకున్న వెంటనే
$+$, $\times$, $<$ల అర్థనిర్దేశాలు కూడా నిర్ణయమైపోతాయి.
$s\colon M\to M\setminus\{z\}$ ప్రమేయం
\tetoken{ఒకటి-ఒకటి}{!!{injective}} మరియు
\tetoken{సంగ్రస్త}{!!{surjective}} రెండూ అయితేనే, అది ప్రామాణిక
నమూనా $\Struct{M}$లో $\Assign{\prime}{M}$గా తగినది. $s$ పరిధిలో $z$ ఉండకూడదు;
లేకపోతే $\lforall[x][\eq/[\Obj 0][x']]$ అసత్యమవుతుంది. ఆ
\tetoken{వాక్యం}{!!{sentence}} $\Struct{N}$లో సత్యం కాబట్టి
$\Struct{M}$ కూడా దాన్ని సత్యం చేయాలి. $s$
\tetoken{ఒకటి-ఒకటి}{!!{injective}} అయి ఉండాలి; ఎందుకంటే
$\Struct{N}$లోని ఉత్తరాధికారి ప్రమేయం $\Assign{\prime}{N}$
\tetoken{ఒకటి-ఒకటి}{!!{injective}}, ఆ $\Assign{\prime}{N}$ ధర్మం
$\Struct{N}$లో సత్యమైన
\tetoken{ఒక వాక్యం}{!!a{sentence}} ద్వారా వ్యక్తమవుతుంది.
ఆ ప్రమేయం \tetoken{సంగ్రస్త}{!!{surjective}} అయి ఉండాలి; లేకపోతే
ఏదో $x\in M\setminus\{z\}$, $s$ పరిధిలో ఉండదు, అంటే
$\lforall[x][(\eq[x][\Obj 0]\lor\lexists[y][\eq[y'][x]])]$ అనే
\tetoken{వాక్యం}{!!{sentence}} $\Struct{M}$లో అసత్యమవుతుంది---కానీ
అది $\Struct{N}$లో సత్యం.
\sourcecorrection{OLTEMODARI-004}{ఇక్కడ సంగ్రస్తత విఫలమైతే
$x$, $s$ యొక్క పరిధిలో ఉండదనాలి; ఆంగ్ల మూలం పొరపాటుగా
“వ్యక్తి క్షేత్రంలో” అంటుంది. ప్రమేయపు వ్యక్తి క్షేత్రం మొత్తం
$M$గానే ఉన్నందున, వాదనకు అవసరమైన పరిధి అనే పదాన్ని వాడాం.}
\end{explain}

\end{document}
